\subsection{Function Quick Reference}
\begingroup
\small
\#include $<$algorithm$>$ \#include $<$numeric$>$ \\
\begin{tabular}{|l|l|p{5.4cm}|} \hline
	\textbf{Algo} & \textbf{Params} &  \textbf{Funcion} \\  \hline
	%swap & e1, e2 &  da vuelta e1,e2 & $1$\\\hline
	sort, stable\_sort & f, l &  ordena el intervalo \\  \hline
	%is\_sorted & f, l &  \textit{bool} si esta ordenado \\  \hline
	nth\_element & f, nth, l & \textit{void} ordena el n-esimo, y \\ && particiona el resto \\  \hline
	fill, fill\_n & f, l / n, elem & \textit{void} llena [f, l) o [f, \\ && f+n) con elem \\  \hline
	lower\_bound, upper\_bound & f, l, elem & \textit{it} al primer / ultimo donde se \\ && puede insertar elem para que\\ && quede ordenada \\  \hline
	binary\_search & f, l, elem & \textit{bool} esta elem en [f, l) \\  \hline
	copy & f, l, resul & hace resul+$i$=f+$i$ $\forall i$ \\  \hline
	find, find\_if, find\_first\_of & f, l, elem & \textit{it} encuentra i $\in$[f,l) tq. i$=$elem, \\ & / pred / f2, l2 & pred(i), i$\in$[f2,l2)\\\hline
	count, count\_if & f, l, elem/pred & cuenta elem, pred(i)\\\hline
	search & f, l, f2, l2 & busca [f2,l2) $\in$ [f,l)\\\hline
	replace, replace\_if & f, l, old & cambia old / pred(i) por new \\ & / pred, new &\\\hline
	reverse & f, l & da vuelta\\\hline
	partition, stable\_partition & f, l, pred & pred(i) ad, !pred(i) atras\\\hline
	%min, max & e1, e2 & men / may & $1$\\\hline
	min\_element, max\_element & f, l, [comp] & \textit{it} min, max de [f,l)\\\hline
	lexicographical\_compare & f1,l1,f2,l2 & \textit{bool} con [f1,l1]<[f2,l2]\\\hline
	next/prev\_permutation & f,l & deja en [f,l) la perm sig, ant\\\hline
	set\_intersection, & f1, l1, f2, l2, res & [res, $\ldots$) la op. de conj\\
	set\_difference, set\_union, & & \\
	set\_symmetric\_difference, & &\\\hline
	push\_heap, pop\_heap, & f, l, e / e / & mete/saca e en heap [f,l), \\
	make\_heap & & hace un heap de [f,l)\\\hline
	is\_heap & f,l & \textit{bool} es [f,l) un heap\\\hline
	accumulate & f,l,i,[op] & \textit{T} $=$ $\sum$/oper de [f,l)\\\hline
	inner\_product & f1, l1, f2, i & \textit{T} $=$ i $+$ [f1, l1) . [f2, $\ldots$ )\\\hline
	partial\_sum & f, l, r, [op] & r+i = $\sum$/oper de [f,f+i] $\forall i \in$[f,l)\\\hline
	%power & e, i, op & \textit{T} = $e^{n}$\\\hline
	\_\_builtin\_ffs& unsigned int & Pos. del primer 1 desde la derecha\\\hline
	\_\_builtin\_clz & unsigned int & Cant. de ceros desde la izquierda.\\\hline
	\_\_builtin\_ctz & unsigned int & Cant. de ceros desde la derecha.\\\hline
	\_\_builtin\_popcount & unsigned int & Cant. de 1’s en x.\\\hline
	\_\_builtin\_parity & unsigned int & 1 si la cantidad de bits en 1 es impar.\\\hline
	\_\_builtin\_XXXXXXll & unsigned ll & = pero para long long's.\\\hline
\end{tabular}
\endgroup
\newpage


\centering
\begingroup
\raggedright
\subsection{C++17 Function Examples}
Runnable examples with expected results in comments. Iterator ranges
use $[first,last)$: the end iterator is excluded. For an inclusive
vector interval $[l,r]$, pass \texttt{begin()+l} and
\texttt{begin()+r+1}. All indices below are zero-based.

\begin{itemize}
\item Reordering: reverse, sort, stable sort, rotate, partial sort,
selection, and permutations.
\item Searching: bounds, binary search, find, counts, predicates,
minimum/maximum, and clamp.
\item Editing: deduplication, erase-remove, replacement, fill,
transform, and filtered copying.
\item Numeric and string helpers: sums, prefixes, differences, dot
products, gcd/lcm, substrings, parsing, and conversion.
\item Sorted set operations and heap algorithms.
\end{itemize}
Each group is a separate function; \texttt{main} runs all examples.
\lstinputlisting[style=C++,showstringspaces=false]{../compile/STLExamples.cpp}
\endgroup


\begingroup
\raggedright
\subsection{Comparators}
For \texttt{sort} and ordered containers, \texttt{comp(a,b)} is true
when \texttt{a} comes before \texttt{b}. In a
\texttt{priority\_queue}, true means \texttt{a} has lower priority
than \texttt{b}; the queue exposes the last element in that order.

\begin{center}
\begin{tabular}{lll}
\textbf{Comparator} & \textbf{\texttt{sort} order} & \textbf{Queue pop order} \\
\hline
\texttt{less<int>} & Smallest first & Largest first \\
\texttt{greater<int>} & Largest first & Smallest first \\
\end{tabular}
\end{center}
Queue output below means repeated \texttt{top()} and \texttt{pop()};
a heap is not a fully sorted array. To make a queue pop in a desired
\texttt{before(a,b)} order, compare \texttt{before(b,a)}.

A comparator must define a strict weak ordering: \texttt{comp(x,x)}
is false; comparisons and equivalence must be transitive. Use strict
comparisons, never \texttt{<=} or \texttt{>=}. Reverse by swapping
arguments, not by negating the result. Keep comparison keys unchanged
while elements are stored in an ordered container or heap.

The runnable examples show ascending/descending integers, pairs with
mixed field directions, custom structs with explicit tie-breaks,
named functions, C++17 lambdas, stable sorting, and set equivalence.
Every example includes its exact output order in comments.
\lstinputlisting[style=C++,showstringspaces=false]{../compile/ComparatorExamples.cpp}


\medskip
\textbf{Additional comparator snippets.}
For the first two lines, assume an \texttt{Edge} type with member
\texttt{w} and an existing vector \texttt{a}. Each declaration below
is a separate snippet.
\cppfile{../more_topics/CustomComparators.cpp}
\endgroup
